\documentclass{beamer}

% You can also use a 16:9 aspect ratio:
%\documentclass[aspectratio=169]{beamer}
\usetheme{TACC16}
\usepackage{graphicx}

% It's possible to move the footer to the right:
%\usetheme[rightfooter]{TACC16}


% Detailed knowledge of application workload characteristics can
% optimize performance of current and future systems. This may sound
% daunting, with many HPC data centers hosting over 2,000 users running
% thousands of applications and millions of jobs per month.  XALT is an
% open source tool developed at the Texas Advanced Computing Center
% (TACC) that collects system usage information to quantitatively report
% how users are using your system. This session will explore the
% benefits of detailed application workload profiling and how the XALT
% tool has helped leading supercomputing sites unlock the power of their
% application usage data. 


\begin{document}
\title[XALT]{XALT: Job-Level Usage Data on Today's Supercomputers.}
\author{Robert McLay} 
\date{Sept. 21, 2022}

% page 1
\frame{\titlepage} 

% page 2
\begin{frame}{XALT: Outline}
  \center{\includegraphics[width=.9\textwidth]{XALT_Sticker.png}}
  \begin{itemize}
    \item What is XALT and what it is not?
    \item How it works: Three Parts
    \item Memory allocation issues
    \item Container issues
    \item Conclusions
  \end{itemize}
\end{frame}

% page 3
\begin{frame}{Understanding what your users are doing}
  \begin{itemize}
    \item Current Version: XALT 2.10.42
    \item What programs, libraries are your users using?
    \item What imports from R, MATLAB, Python?
    \item What are the top programs by core-hours? by counts? by users?
    \item System, User or Built by Other executables?
    \item Are Executables implemented in C/C++/Fortran/Others?
    \item Track MPI task and/or Threading (\$OMP\_NUMTHREADS)
    \item Function Tracking
    \item Census Taker, Not a performance tool!
  \end{itemize}
\end{frame}

% page 4
\begin{frame}{How do TACC and other sites use XALT data?}
  \begin{itemize}
    \item Some sites feed XALT data into an ELK stack to visualize data
    \item We use a MySQL db to store  the data and analyze weekly.
    \item Sites could do both ELK and MySQL DB (or some other DB)
    \item TACC generates a weekly report on usage (last 30, last 365) days
    \item We design our benchmark codes on actual usage not guesses
    \item Usage of gromacs, amber, ...?
  \end{itemize}
\end{frame}

% page 4
\begin{frame}{xalt\_usage\_report.py}
  \begin{itemize}
    \item Overall Job counts: system (5\%) user (95\%)
    \item It reports Top Execs by Core-Hours, Number of Runs,
      Number of users for All, MPI Only, Scalar
    \item Top Module usage
    \item Compiler usage
    \item Libraries from modules Usage
  \end{itemize}
\end{frame}

% page 4
\begin{frame}{XALT Monthly Zoom Mtg}
  \begin{itemize}
    \item Previous Topics:
      \begin{itemize}
        \item Controlling which executable get tracked
        \item Using Lmod to build the reverse map from directories to modules
        \item Using \$XALT\_TRACING to debug XALT
        \item How to test XALT at your site
        \item Examples of what data a site can report on from the XALT DB.
      \end{itemize}
    \item Future Topics:
      \begin{itemize}
        \item Tracking Packages: Python, R, MATLAB
      \end{itemize}
    \item Zoom Mtg Usually 3rd Thurday at 15:30 UTC (9:30 US Central)
    \item See Mailing list or https://github.com/xalt/xalt/wiki for
      details
    \item Next Meeting Feb 17 at 16:00 UTC (10:00 US Central)
  \end{itemize}
\end{frame}

% page 4
\begin{frame}{Design Goals}
  \begin{itemize}
    \item Be extremely light-weight
    \item How many use a library or application?
    \item What functions are users calling in system modules
    \item Collect Data into a Database for analysis.
  \end{itemize}
\end{frame}

% page 5
\begin{frame}{How does XALT work?}
  \begin{itemize}
    \item LD Wrapper
    \item ELF Trick to track execution
    \item Generate Json records
    \item Transport to DB
    \item Analyze database.
  \end{itemize}
\end{frame}

% page 6
\begin{frame}{Design: LD Wrapper}
  \begin{itemize}
    \item XALT wraps the linker to enable tracking of exec's
    \item The linker (ld) wrapper intercepts the user link line.
    \item Generate assembly code: key-value pairs
    \item Capture tracemap output from ld
    \item Transmit collected data in *.json format
    \item Optionally add codes that executes before main() and after
      main() completes for static builds.
    \item Having the LD Wrapper is helpful but not required for XALT
      to work.
  \end{itemize}
\end{frame}


% page 7
\begin{frame}{Elf Trick (I)}
  \begin{itemize}
    \item ELF is the binary format for Linux
    \item ELF has many hooks
    \item XALT uses two hooks to run before/after main()
  \end{itemize}
\end{frame}

% page 8
\begin{frame}[fragile]
    \frametitle{ELF Trick (II)}
 {\small
    \begin{semiverbatim}
{\color{blue}hello.c}:
#include <stdio.h>
int main() \{ printf("Hello World!\\n"): return 0; \}

{\color{blue}xalt.c}:
#include <stdio.h>
void myinit(int argc, char **argv)
\{ fprintf(stderr, "This is run before main()\textbackslash{}n"); \}
void myfini()
\{ fprintf(stderr,"This is run after main()\textbackslash{}n");\}
static \_\_attribute\_\_((section(".init_array")))
     typeof(myinit) *\_\_init = myinit;
static \_\_attribute\_\_((section(".fini_array")))
     typeof(myfini) *\_\_fini = myfini;
    \end{semiverbatim}
}
\end{frame}

% page 9
\begin{frame}[fragile]
    \frametitle{ELF Trick (III)}
 {\small
    \begin{semiverbatim}
{\color{red}\texttt{\% ./hello}}
Hello World!

{\color{red}\texttt{\% gcc -c -fPIC xalt.c;}}
{\color{red}\texttt{\% gcc -o libxalt.so -fPIC -shared xalt.o}}
{\color{red}\texttt{\% LD\_PRELOAD=./libxalt.so ./hello}}
This is run before main()
Hello World!
This is run after main()
    \end{semiverbatim}
}
\end{frame}

% page 11
\begin{frame}{Transmission to DB}
  \begin{itemize}
    \item File: collect nightly/hourly/...
    \item Syslog: Use Syslog filtering (or ELK)
    \item Curl: send directly (typically for ELK)
  \end{itemize}
\end{frame}

% page 12
\begin{frame}{Lmod to XALT connection (I)}
  \begin{itemize}
    \item Optional support to connect paths to modules
    \item Lmod spider walks entire module tree.
    \item Can build a reverse map from paths to modules
    \item Can map program \& libraries to modules.
    \item /opt/apps/i15/mv2\_2\_1/phdf5/1.8.14/lib/libhdf5.so.9
      $\Rightarrow$ phdf5/1.8.14(intel/15.02:mvapich2/2.1)
    \item Tmod Sites can still use Lmod to build the reverse map.
  \end{itemize}
\end{frame}

% page 13
\begin{frame}{Installing XALT}
  \begin{itemize}
    \item Easy: \texttt{./configure ...; make install}
    \item Harder: Reverse Map from Lmod?
    \item Harder: Site config.py file
    \item Harder: Setup Transport Json records
    \item Harder: Setup VM to hold database
    \item Or: Setup your own way to handle the data
  \end{itemize}
\end{frame}

% page 14
\begin{frame}{Site config.py (I)}
  \begin{itemize}
    \item Each site must configure to match their setup
    \item Compute node names?
    \item What executables to track or ignore?
    \item What python packages to track or ignore?
    \item What sampling rules to use?
  \end{itemize}
\end{frame}

% page 15
\begin{frame}{Site config.py (II)}
  \begin{itemize}
    \item XALT use config.py to create *.h, *.lex *.py files during build.
    \item Provides xalt\_configuration\_report C++ program to know how configured.
    \item Config.py file only used when building XALT.
    \item Any changes to Config.py require a re-install of XALT.
  \end{itemize}
\end{frame}

% page 16
\begin{frame}{Hostname, Path and Env Filtering}
  \begin{itemize}
    \item Uses FLEX to compile in patterns
    \item Use regex expression to control what to keep and ignore.
    \item These Flex $\Rightarrow$ C routine provide fast regex parsing.
  \end{itemize}
\end{frame}

% page 17
\begin{frame}[fragile]
    \frametitle{TACC\_config.py}
 {\small
    \begin{semiverbatim}
hostname\_patterns = [
  ['KEEP', '^c[0-9][0-9][0-9]-[0-9][0-9][0-9]\textbackslash{}..*']
  ['KEEP', '^nid[0-9][0-9][0-9][0-9][0-9].*']
]
path\_patterns = [
  ['PKGS',  r'.*\textbackslash{}/python[0-9.]*'],
  ['PKGS',  r'.*\textbackslash{}/R'],
  ['KEEP',  r'^/\textbackslash{}usr\textbackslash{}/bin\textbackslash{}/ddt'],
  ['SKIP',  r'^\textbackslash{}/bin\textbackslash{}/.*'],
]
env\_patterns = [
  [ 'SKIP', r'^MKLROOT=.*' ],
  [ 'SKIP', r'^MKL\_DIR=.*' ],
  [ 'KEEP', r'^I\_MPI\_INFO\_NUMA\_NODE\_NUM=.*'],
]
    \end{semiverbatim}
}
\end{frame}

% page 18
\begin{frame}{How sampling works}
  \begin{itemize}
    \item Changed design to deal with the overload of XALT data
    \item Only generate records if plan to save.
    \item All Non-mpi executions only produce end records.
    \item Small MPI execution sample, Large MPI executions record.
  \end{itemize}
\end{frame}

% page 19
\begin{frame}{Sampling Non-MPI programs}
  \begin{itemize}
    \item XALT has sampling rules (site configurable!)
    \item TACC rules are:
    \item 0 mins $<$ 30 mins $\Rightarrow$ 0.01\% recorded  
    \item 30 mins $<$ 120 mins $\Rightarrow$ 1\% recorded  
    \item 120 mins $<  \infty \Rightarrow$ 100\% recorded
    \item Can now track/sample perl, awk, sed, gzip etc
  \end{itemize}
\end{frame}

% page 20
\begin{frame}{Sampling MPI programs}
  \begin{itemize}
    \item Some users are using many short MPI programs to train Deep
      Learning engine
    \item TACC rules are:
    \item Task counts $<$ 128 tasks are sampled
      \begin{itemize}
        \item 0 mins $<$ 15 mins $\Rightarrow$ 0.01\% recorded  
        \item 15 mins $<$ 30 mins $\Rightarrow$ 1\% recorded  
        \item 30 mins $<  \infty \Rightarrow$ 100\% recorded
      \end{itemize}
    \item Task counts $>=$ 128 task are always recorded independent of runtime.
    \item Need to Capture long running MPI progs that never end.
  \end{itemize}
\end{frame}

% page 21
\begin{frame}{XALT is now linking with everything!}
  \begin{itemize}
    \item I feel like a developer on every team.
    \item XALT shares name space -> obfuscation
    \item XALT shares memory allocation
    \item Containers: XALT cannot depend on all system libs.
  \end{itemize}
\end{frame}

% page 23
\begin{frame}{Memory Collusion w/ user code }
  \begin{itemize}
    \item Not all user programs allocate and free memory correctly (Surprise!)
    \item XALT would sometimes fail when freeing memory after main()
    \item Result: XALT allocates but doesn't free memory after main()
  \end{itemize}
\end{frame}

% page 24
\begin{frame}{Containers}
  \begin{itemize}
    \item XALT requires libuuid.so on host
    \item Not on Containers
    \item XALT has to use dlopen()/dlsym() to use libuuid.so
  \end{itemize}
\end{frame}

% page 25
\begin{frame}{Containers (II)}
  \begin{itemize}
    \item XALT copies system libuuid.so to XALT install dir.
    \item XALT does: \texttt{dlopen("XALT\_INSTALL\_DIR/lib64/libuuid.so")}
    \item XALT then uses dlsym() to connect to libuuid routines
  \end{itemize}
\end{frame}

% page 28
\begin{frame}{Tracking Python packages}
  \begin{itemize}
    \item Help from Riccardo Murri
    \item sitecustomize.py
    \item It is run by any Python if found.
    \item All Pythons uses sys.meta\_path to locate files to import
    \item Can register object to capture imports.
    \item Just add location to PYTHONPATH
  \end{itemize}
\end{frame}

% page 29
\begin{frame}[fragile]
    \frametitle{Filtering python packages}
 {\small
    \begin{semiverbatim}
\{ 'k\_s':'SKIP','kind':'path','patt': r"^[^/]" \},
\{ 'k\_s':'SKIP','kind':'name','patt': r"^_"    \},
\{ 'k\_s':'SKIP','kind':'name','patt': r".*\textbackslash." \},
\{ 'k\_s':'KEEP','kind':'path','patt': r".*/.local/" \},
    \end{semiverbatim}
}
\end{frame}

% page 30
\begin{frame}{New program: xalt\_extract\_record}
  \begin{itemize}
    \item This program reads the watermark.
    \item Find out who built this program on what machine
    \item Find out what modules where used.
    \item Where was it build. 
  \end{itemize}
\end{frame}

% page 31
\begin{frame}[fragile]
    \frametitle{Example of xalt\_extract\_record output}
 {\tiny
    \begin{semiverbatim}
****************************************
XALT Watermark: hello
****************************************
Build\_CWD            /home/user/t/hello
Build\_Epoch          1510257139.4624
Build\_LMFILES        /apps/mfiles/intel/17.0.4.lua:...
Build\_LOADEDMODULES  intel/18.0.4:impi/18.0.3:TACC:...
Build\_OS             Linux 3.10.0-514.26.2.el7.x86_64
Build\_Syshost        stampede2
Build\_UUID           586d5943-67eb-480b-a2fe-35e87a1f22c7
Build\_User           mclay
Build\_compiler       icc
Build\_date           Fri Jun 09 13:52:19 2019
Build\_host           c455-011.stampede2.tacc.utexas.edu
XALT\_Version         2.7
    \end{semiverbatim}
}
\end{frame}

% page 33
\begin{frame}{New Feature: Track Singularity Container Usage}
  \begin{itemize}
    \item Sites can configure their Singularity script to include XALT
    \item It works well with syslog, file  or curl transfer of data
    \item Thanks to Scott McMillan from NVIDIA for the contribution.
  \end{itemize}
\end{frame}

% page 34
\begin{frame}[fragile]
    \frametitle{Debugging XALT}
 {\tiny
    \begin{semiverbatim}
\$ XALT\_TRACING=yes XALT\_SAMPLING=no ./hello

myinit(0/1,LD_PRELOAD,/path/to/hello)\{
    -> Setting up signals
    -> Leaving myinit
\}

Hello World!

myfini(0/1,LD_PRELOAD,/path/to/hello)\{
    -> XALT_SAMPLING = "no" All programs tracked!
  Recording State at end of scalar user program:
    XALT_EXECUTABLE_TRACKING=no /path/to/xalt_run_submission ...

  xalt_run_submission(.zzz) \{
    building UUID: ...
    Extracted recordT from executable
    Built userT, userDT, scheduler: SLURM
    Using XALT_TRANSMISSION_STYLE: file
    cmdlineA: ["./hello"]
    Built json string
  \}
    -> leaving myfini
\}
    \end{semiverbatim}
}
\end{frame}


% page 
\begin{frame}{Internal Changes}
  \begin{itemize}
    \item Use of CRC to fix dup UUIDs problem
    \item Pre-Ingestion Filter
    \item Better Support for compilers: rustc, chpl, ocaml, ...
    \item Better support for mpi wrapper compilers: mpicc(icc), ...
  \end{itemize}
\end{frame}


% page 36
\begin{frame}{Conclusion}
  \center{\includegraphics[width=.9\textwidth]{XALT_Sticker.png}}
  \begin{itemize}
    \item Lmod:
      \begin{itemize}
        \item Source: github.com/TACC/lmod.git, lmod.sf.net
        \item Documentation: lmod.readthedocs.org
      \end{itemize}
    \item XALT:
      \begin{itemize}
        \item Source: github.com/xalt/xalt.git, xalt.sf.net
        \item Documentation: XALT 2 $\Rightarrow$ xalt.readthedocs.org
        \item Join mailing list: https://sourceforge.net/projects/xalt/lists/xalt-users
      \end{itemize}
  \end{itemize}
\end{frame}

%
\begin{frame}{License}
\footnotesize
\textcopyright The University of Texas at Austin, \the\year\\
\vspace{0.5 cm}
This work is licensed under the Creative Commons Attribution Non-Commercial 3.0 Unported License. To view a copy of this license, 
visit http://creativecommons.org/licenses/by-nc/3.0/\\
\vspace{0.5 cm}
When attributing this work, please use the following text: 
"Designing and Administering Large-scale Systems", Texas Advanced Computing Center, \the\year. Available under a Creative Commons Attribution Non-Commercial 3.0 Unported License.
\newline
\newline
\includegraphics[width=1cm]{./figures/cc.png}
\end{frame}

\end{document}
